\subsection{A KL-type control of $\Lambda_{\mathrm{var}}$ for the optimistic kernel}

In this subsection we revisit the quantity
\begin{align*}
    \Lambda_{\mathrm{var}, q}(E, p; f)
     :=
    \E\left[
        D_{H^\star}\left(
        H'(q) - \eta E(X, f(X)), H'(q)
        \right)
        \right],
\end{align*}
where $X \sim p$.
Our goal is twofold:
\begin{enumerate}
    \item to show that this quantity is already of KL type, and
    \item for the optimistic kernel at the end of the previous subsection, to bound it by an ordinary KL mutual information rather than by a $\chi^2$-type quantity.
\end{enumerate}

Throughout this subsection we assume the following.
\begin{itemize}
    \item $q$ is a probability density on $\K$ such that $q(t) > 0$ for almost every $t \in \K$.
    \item Whenever we use differential entropy below, we assume it is finite. This is a standard regularity assumption and is satisfied, for example, under the usual full-dimensional log-concavity assumptions with mild boundary behavior.
    \item For the optimistic-kernel bound we assume $0 \leq f(x) \leq 1$ for all $x \in \K$ and we set $b_x(t) = 0$ as in the last display of the previous subsection.
\end{itemize}

\paragraph{Notation.}
For readability, when $f$ is fixed we write
\begin{align}
    g_x(t) := E(x, f(x))(t).
    \label{eq:appendix_gx_def}
\end{align}
Thus the general definition becomes
\begin{align}
    \Lambda_{\mathrm{var}, q}(E, p; f)
    =
    \E\left[
        D_{H^\star}\left(
        H'(q) - \eta g_X, H'(q)
        \right)
    \right].
    \label{eq:appendix_lambda_var_general_def}
\end{align}

\begin{proposition}[The CEXP variance term is a generalized KL divergence]
\label{prop:appendix_generalized_kl_form}
For each $x \in \K$, define
\begin{align}
    \widehat q_x(t) := q(t)e^{-\eta g_x(t)}.
    \label{eq:appendix_qhat_def}
\end{align}
Then
\begin{align}
    \Lambda_{\mathrm{var}, q}(E, p; f)
    =
    \E\left[D_H\left(q, \widehat q_X\right)\right],
    \label{eq:appendix_generalized_kl_identity}
\end{align}
where
\begin{align}
    D_H(u,v) = \int_\K \left(u(t)\log\frac{u(t)}{v(t)} - u(t) + v(t)\right)dt.
    \label{eq:appendix_generalized_kl_formula}
\end{align}
Equivalently,
\begin{align}
    \Lambda_{\mathrm{var}, q}(E, p; f)
    =
    \int_\K\int_\K q(t)\left(e^{-\eta g_x(t)} + \eta g_x(t) - 1\right)p(x)dt\,dx.
    \label{eq:appendix_lambda_var_integral_formula}
\end{align}
\end{proposition}

\begin{proof}
We expand the Bregman divergence associated with $H^\star$.
By definition,
\begin{align*}
     & D_{H^\star}\left(H'(q)-\eta g_x, H'(q)\right) \\
     & = H^\star\left(H'(q)-\eta g_x\right) - H^\star\left(H'(q)\right) + \eta \langle g_x, q\rangle.
\end{align*}
Recall that
\begin{align*}
    H^\star(v) = \int_\K e^{v(t)}dt,
    \qquad
    H'(q) = \log q.
\end{align*}
Therefore,
\begin{align}
    H^\star\left(H'(q)\right)
     & = \int_\K e^{\log q(t)}dt
    = \int_\K q(t)dt
    = 1,
    \label{eq:appendix_hstar_at_hprime_q}
\end{align}
and also
\begin{align}
    H^\star\left(H'(q)-\eta g_x\right)
     & = \int_\K e^{\log q(t)-\eta g_x(t)}dt
    = \int_\K q(t)e^{-\eta g_x(t)}dt.
    \label{eq:appendix_hstar_shifted}
\end{align}
Plugging \eqref{eq:appendix_hstar_at_hprime_q} and \eqref{eq:appendix_hstar_shifted} into the definition of the Bregman divergence gives
\begin{align}
     & D_{H^\star}\left(H'(q)-\eta g_x, H'(q)\right) \nonumber \\
     & = \int_\K q(t)e^{-\eta g_x(t)}dt - 1 + \eta \int_\K g_x(t)q(t)dt \nonumber \\
     & = \int_\K q(t)\left(e^{-\eta g_x(t)} + \eta g_x(t) - 1\right)dt.
    \label{eq:appendix_expand_hstar_bregman}
\end{align}
Now compare this with $D_H(q,\widehat q_x)$, where $\widehat q_x$ is defined in \eqref{eq:appendix_qhat_def}:
\begin{align*}
    D_H(q,\widehat q_x)
     & = \int_\K \left(q(t)\log\frac{q(t)}{\widehat q_x(t)} - q(t) + \widehat q_x(t)\right)dt.
\end{align*}
Since $\widehat q_x(t)=q(t)e^{-\eta g_x(t)}$, we have
\begin{align*}
    \log\frac{q(t)}{\widehat q_x(t)}
    = \log\left(e^{\eta g_x(t)}\right)
    = \eta g_x(t).
\end{align*}
Hence
\begin{align*}
    D_H(q,\widehat q_x)
     & = \int_\K \left(q(t)\eta g_x(t) - q(t) + q(t)e^{-\eta g_x(t)}\right)dt \\
     & = \int_\K q(t)\left(e^{-\eta g_x(t)} + \eta g_x(t) - 1\right)dt.
\end{align*}
This agrees exactly with \eqref{eq:appendix_expand_hstar_bregman}. Averaging over $X \sim p$ proves
\eqref{eq:appendix_generalized_kl_identity}, and writing out the expectation yields
\eqref{eq:appendix_lambda_var_integral_formula}.
\end{proof}

The previous proposition already shows that $\Lambda_{\mathrm{var}}$ is a KL-type quantity, but the second argument of $D_H$ is not normalized in general.
The next proposition isolates the ordinary KL part explicitly.

\begin{proposition}[Ordinary KL plus a mass correction]
\label{prop:appendix_normalized_kl_decomposition}
Let $\widehat q_x$ be as in \eqref{eq:appendix_qhat_def}, and define
\begin{align}
    Z_x := \int_\K \widehat q_x(t)dt,
    \qquad
    \overline q_x(t) := \frac{\widehat q_x(t)}{Z_x}.
    \label{eq:appendix_normalization_def}
\end{align}
Then $\overline q_x$ is a probability density and
\begin{align}
    D_H(q,\widehat q_x)
    =
    \int_\K q(t)\log\frac{q(t)}{\overline q_x(t)}dt
    + Z_x - 1 - \log Z_x.
    \label{eq:appendix_kl_plus_mass_identity}
\end{align}
Consequently,
\begin{align}
    \Lambda_{\mathrm{var}, q}(E, p; f)
    =
    \E\left[\int_\K q(t)\log\frac{q(t)}{\overline q_X(t)}dt\right]
    + \E\left[Z_X - 1 - \log Z_X\right].
    \label{eq:appendix_lambda_var_kl_plus_mass}
\end{align}
\end{proposition}

\begin{proof}
Since $\widehat q_x = Z_x\overline q_x$, we have
\begin{align*}
    \log\frac{q(t)}{\widehat q_x(t)}
    = \log\frac{q(t)}{\overline q_x(t)} - \log Z_x.
\end{align*}
Substitute this into the formula for $D_H(q,\widehat q_x)$:
\begin{align*}
    D_H(q,\widehat q_x)
     & = \int_\K q(t)\log\frac{q(t)}{\overline q_x(t)}dt - \int_\K q(t)\log Z_x\,dt - \int_\K q(t)dt + \int_\K \widehat q_x(t)dt.
\end{align*}
Now each of the last three terms is elementary:
\begin{align*}
    \int_\K q(t)\log Z_x\,dt
     & = \log Z_x \int_\K q(t)dt
    = \log Z_x, \\
    \int_\K q(t)dt
     & = 1, \\
    \int_\K \widehat q_x(t)dt
     & = Z_x.
\end{align*}
Therefore,
\begin{align*}
    D_H(q,\widehat q_x)
     & = \int_\K q(t)\log\frac{q(t)}{\overline q_x(t)}dt + Z_x - 1 - \log Z_x,
\end{align*}
which is exactly \eqref{eq:appendix_kl_plus_mass_identity}. Averaging over $X$ yields
\eqref{eq:appendix_lambda_var_kl_plus_mass}.
\end{proof}

Thus, even before using the optimistic kernel, the CEXP variance term is already of the form
\begin{align*}
    \text{ordinary KL divergence in the $t$-variable}
    \;+
    \text{a one-dimensional mass correction}.
\end{align*}

We now turn to the optimistic kernel from the end of the previous subsection and derive an explicit upper bound.

\begin{proposition}[A useful identity for the optimistic kernel]
\label{prop:appendix_optimistic_kernel_identity}
Let
\begin{align*}
    X = \lambda\sum_{k=0}^\infty (1-\lambda)^k T_k,
\end{align*}
where $T_0,T_1,\dots$ are i.i.d. with common density $q$, and let $p_X$ be the density of $X$.
Define
\begin{align}
    r(x,t) := \frac{p_{X\mid T_0=t}(x)}{p_X(x)}.
    \label{eq:appendix_r_def}
\end{align}
Then the optimistic kernel satisfies
\begin{align}
    s_x^\lambda(t)
    = 1 - \frac{1}{\lambda} + \frac{1}{\lambda}r(x,t).
    \label{eq:appendix_s_in_terms_of_r}
\end{align}
Moreover, if $Q := P_X P_{T_0}$, then
\begin{align}
    \E_Q[r] = 1.
    \label{eq:appendix_r_mean_one}
\end{align}
\end{proposition}

\begin{proof}
From the definition of $X$ we have the decomposition
\begin{align*}
    X
    = \lambda T_0 + (1-\lambda)X',
    \qquad
    X' := \lambda\sum_{k=0}^\infty (1-\lambda)^k T_{k+1}.
\end{align*}
Since $X'$ is independent of $T_0$ and $X' \stackrel{d}{=} X$, conditioning on $T_0=t$ gives
\begin{align*}
    X\mid(T_0=t) = \lambda t + (1-\lambda)X'.
\end{align*}
Therefore,
\begin{align}
    p_{X\mid T_0=t}(x)
    = \frac{1}{(1-\lambda)^d}
    p_X\left(\frac{x-\lambda t}{1-\lambda}\right).
    \label{eq:appendix_conditional_density}
\end{align}
Comparing this with the definition of $s_x^\lambda(t)$ from the previous subsection yields
\eqref{eq:appendix_s_in_terms_of_r} immediately.

Next we show \eqref{eq:appendix_r_mean_one}. By definition of $r$,
\begin{align*}
    \E_Q[r]
     & = \int_\K\int_\K p_X(x)q(t)\frac{p_{X\mid T_0=t}(x)}{p_X(x)}dt\,dx \\
     & = \int_\K q(t)\left(\int_\K p_{X\mid T_0=t}(x)dx\right)dt \\
     & = \int_\K q(t)dt
    = 1.
\end{align*}
\end{proof}

The next step is the delicate one: justifying carefully why the loss $f$ can be replaced by $1$ when $0 \leq f \leq 1$.
The point is that $s_x^\lambda(t)$ may be negative, so this is not a simple monotonicity argument.

\begin{lemma}[The worst case is $f \equiv 1$]
\label{lem:appendix_f_one_reduction}
Assume $0 \leq f(x) \leq 1$ for all $x \in \K$ and $b_x(t)=0$.
Then
\begin{align}
    \Lambda_{\mathrm{var}, q}(E_\lambda, p_X; f)
    \leq
    \int_\K\int_\K \left(e^{-\eta s_x^\lambda(t)} + \eta s_x^\lambda(t) - 1\right)p_X(x)q(t)dt\,dx.
    \label{eq:appendix_lambda_var_reduce_to_one}
\end{align}
\end{lemma}

\begin{proof}
Define
\begin{align}
    \phi_\eta(u) := e^{-\eta u} + \eta u - 1.
    \label{eq:appendix_phi_def}
\end{align}
We claim that for every $a \in [0,1]$ and every $u \in \R$,
\begin{align}
    \phi_\eta(au) \leq \phi_\eta(u).
    \label{eq:appendix_phi_contraction}
\end{align}
Once this is proved, we can apply it with $a=f(x)$ and $u=s_x^\lambda(t)$.

So let us prove \eqref{eq:appendix_phi_contraction} carefully.
First,
\begin{align*}
    \phi_\eta(0) = e^0 + 0 - 1 = 0.
\end{align*}
Second, for every $u \in \R$ we have $\phi_\eta(u) \geq 0$, because the elementary inequality $e^z \geq 1+z$ applied with $z=-\eta u$ gives
\begin{align*}
    e^{-\eta u} \geq 1-\eta u
    \qquad \Longrightarrow \qquad
    e^{-\eta u} + \eta u - 1 \geq 0.
\end{align*}
Third,
\begin{align*}
    \phi_\eta''(u) = \eta^2 e^{-\eta u} > 0,
\end{align*}
so $\phi_\eta$ is convex on all of $\R$.
Now, since $au = a\cdot u + (1-a)\cdot 0$, convexity gives
\begin{align*}
    \phi_\eta(au)
     & = \phi_\eta\left(a u + (1-a)0\right) \\
     & \leq a\phi_\eta(u) + (1-a)\phi_\eta(0) \\
     & = a\phi_\eta(u) \\
     & \leq \phi_\eta(u),
\end{align*}
because $0 \leq a \leq 1$ and $\phi_\eta(u) \geq 0$.
This proves \eqref{eq:appendix_phi_contraction}.

Since $b_x(t)=0$, we have $g_x(t)=f(x)s_x^\lambda(t)$. Therefore, by
\eqref{eq:appendix_lambda_var_integral_formula},
\begin{align*}
    \Lambda_{\mathrm{var}, q}(E_\lambda, p_X; f)
     & = \int_\K\int_\K q(t)\phi_\eta\left(f(x)s_x^\lambda(t)\right)p_X(x)dt\,dx \\
     & \leq \int_\K\int_\K q(t)\phi_\eta\left(s_x^\lambda(t)\right)p_X(x)dt\,dx,
\end{align*}
which is exactly \eqref{eq:appendix_lambda_var_reduce_to_one}.
\end{proof}

Thus, for upper bounds, it suffices to understand the case $f \equiv 1$.
We now rewrite the right-hand side of \eqref{eq:appendix_lambda_var_reduce_to_one} as an $f$-divergence between the joint law and the product law.

\begin{proposition}[Exact $f$-divergence representation]
\label{prop:appendix_f_divergence_representation}
Let
\begin{align*}
    P := P_{X,T_0},
    \qquad
    Q := P_X P_{T_0}.
\end{align*}
Then
\begin{align}
    \frac{dP}{dQ}(x,t) = r(x,t),
    \label{eq:appendix_r_is_likelihood_ratio}
\end{align}
and if we define
\begin{align}
    f_{\eta,\lambda}(u)
    := e^{\eta(\frac{1}{\lambda}-1)}\left(e^{-\frac{\eta}{\lambda}u} - e^{-\frac{\eta}{\lambda}} + \frac{\eta}{\lambda}e^{-\frac{\eta}{\lambda}}(u-1)\right),
    \label{eq:appendix_f_generator_def}
\end{align}
then
\begin{align}
    \Lambda_{\mathrm{var}, q}(E_\lambda, p_X; f)
    \leq
    e^{-\eta} + \eta - 1 + D_{f_{\eta,\lambda}}(P_{X,T_0}\,\|\,P_XP_{T_0}),
    \label{eq:appendix_f_divergence_bound}
\end{align}
where
\begin{align*}
    D_{f_{\eta,\lambda}}(P\|Q)
    := \E_Q\left[f_{\eta,\lambda}\left(\frac{dP}{dQ}\right)\right].
\end{align*}
\end{proposition}

\begin{proof}
By definition of conditional density,
\begin{align*}
    \frac{dP}{dQ}(x,t)
    = \frac{p_{X\mid T_0=t}(x)}{p_X(x)}
    = r(x,t),
\end{align*}
which proves \eqref{eq:appendix_r_is_likelihood_ratio}.
Now start from the upper bound in \eqref{eq:appendix_lambda_var_reduce_to_one} and use
\eqref{eq:appendix_s_in_terms_of_r}:
\begin{align}
     & \Lambda_{\mathrm{var}, q}(E_\lambda, p_X; f) \nonumber \\
     & \leq \E_Q\left[\phi_\eta\left(1 - \frac{1}{\lambda} + \frac{1}{\lambda}r\right)\right].
    \label{eq:appendix_before_phi_expansion}
\end{align}
We now expand the argument of $\phi_\eta$ carefully. Using the definition
\eqref{eq:appendix_phi_def},
\begin{align}
     & \phi_\eta\left(1 - \frac{1}{\lambda} + \frac{1}{\lambda}r\right) \nonumber \\
     & = \exp\left(-\eta\left(1 - \frac{1}{\lambda} + \frac{1}{\lambda}r\right)\right)
    + \eta\left(1 - \frac{1}{\lambda} + \frac{1}{\lambda}r\right) - 1 \nonumber \\
     & = e^{\eta(\frac{1}{\lambda}-1)} e^{-(\eta/\lambda)r}
    + \eta - 1 + \frac{\eta}{\lambda}(r-1).
    \label{eq:appendix_phi_expanded}
\end{align}
Taking expectation with respect to $Q$ and using \eqref{eq:appendix_r_mean_one} gives
\begin{align}
     & \E_Q\left[\phi_\eta\left(1 - \frac{1}{\lambda} + \frac{1}{\lambda}r\right)\right] \nonumber \\
     & = \eta - 1 + e^{\eta(\frac{1}{\lambda}-1)}\E_Q\left[e^{-(\eta/\lambda)r}\right].
    \label{eq:appendix_phi_after_mean_one}
\end{align}
Now define $f_{\eta,\lambda}$ by \eqref{eq:appendix_f_generator_def}. Then
\begin{align*}
     & \E_Q\left[f_{\eta,\lambda}(r)\right] \\
     & = e^{\eta(\frac{1}{\lambda}-1)}\left(\E_Q\left[e^{-(\eta/\lambda)r}\right] - e^{-\eta/\lambda} + \frac{\eta}{\lambda}e^{-\eta/\lambda}\E_Q[r-1]\right).
\end{align*}
Since $\E_Q[r-1]=0$, this becomes
\begin{align*}
    \E_Q\left[f_{\eta,\lambda}(r)\right]
    = e^{\eta(\frac{1}{\lambda}-1)}\left(\E_Q\left[e^{-(\eta/\lambda)r}\right] - e^{-\eta/\lambda}\right).
\end{align*}
Finally,
\begin{align*}
    e^{\eta(\frac{1}{\lambda}-1)}e^{-\eta/\lambda}
    = e^{-\eta},
\end{align*}
so \eqref{eq:appendix_phi_after_mean_one} can be rewritten as
\begin{align*}
    \E_Q\left[\phi_\eta\left(1 - \frac{1}{\lambda} + \frac{1}{\lambda}r\right)\right]
    = e^{-\eta} + \eta - 1 + \E_Q\left[f_{\eta,\lambda}(r)\right].
\end{align*}
Since $r = dP/dQ$, the last expectation is exactly
$D_{f_{\eta,\lambda}}(P_{X,T_0}\,\|\,P_XP_{T_0})$. Combining this with
\eqref{eq:appendix_before_phi_expansion} proves \eqref{eq:appendix_f_divergence_bound}.
\end{proof}

The previous proposition shows that the optimistic-kernel variance term is controlled not by a quadratic generator, but by a much softer $f$-divergence.
The next step is to compare that $f$-divergence with ordinary KL.

\begin{proposition}[Comparison with ordinary KL]
\label{prop:appendix_compare_f_to_kl}
Let
\begin{align}
    h(u) := u\log u - u + 1,
    \label{eq:appendix_kl_generator}
\end{align}
which is the standard generator of the KL divergence.
Then for every $u > 0$,
\begin{align}
    f_{\eta,\lambda}(u)
    \leq
    \frac{\eta}{e\lambda}e^{\eta(\frac{1}{\lambda}-1)} h(u).
    \label{eq:appendix_pointwise_f_vs_kl}
\end{align}
Consequently,
\begin{align}
    D_{f_{\eta,\lambda}}(P_{X,T_0}\,\|\,P_XP_{T_0})
    \leq
    \frac{\eta}{e\lambda}e^{\eta(\frac{1}{\lambda}-1)}
    \mathrm{KL}(P_{X,T_0}\,\|\,P_XP_{T_0}).
    \label{eq:appendix_f_vs_kl_divergence}
\end{align}
\end{proposition}

\begin{proof}
We compare the two generators by comparing their second derivatives.
First,
\begin{align*}
    h''(u) = \frac{1}{u},
    \qquad u>0.
\end{align*}
Next, differentiating \eqref{eq:appendix_f_generator_def} twice gives
\begin{align*}
    f_{\eta,\lambda}''(u)
    = \left(\frac{\eta}{\lambda}\right)^2 e^{\eta(\frac{1}{\lambda}-1)} e^{-(\eta/\lambda)u}.
\end{align*}
Therefore,
\begin{align}
    \frac{f_{\eta,\lambda}''(u)}{h''(u)}
    = \left(\frac{\eta}{\lambda}\right)^2 e^{\eta(\frac{1}{\lambda}-1)} u e^{-(\eta/\lambda)u}.
    \label{eq:appendix_second_derivative_ratio}
\end{align}
The only nontrivial part is to bound the factor $u e^{-(\eta/\lambda)u}$ uniformly in $u$.
Set
\begin{align*}
    a := \frac{\eta}{\lambda} > 0.
\end{align*}
Then the function $u \mapsto ue^{-au}$ has derivative
\begin{align*}
    \frac{d}{du}\left(ue^{-au}\right) = e^{-au}(1-au),
\end{align*}
which vanishes exactly at $u = 1/a$.
Thus the maximum is attained at $u = 1/a$, and
\begin{align*}
    \sup_{u>0} ue^{-au}
    = \frac{1}{ae}.
\end{align*}
Substituting this bound into \eqref{eq:appendix_second_derivative_ratio} gives
\begin{align}
    f_{\eta,\lambda}''(u)
    \leq
    \frac{\eta}{e\lambda}e^{\eta(\frac{1}{\lambda}-1)} h''(u)
    \qquad \text{for all } u>0.
    \label{eq:appendix_second_derivative_domination}
\end{align}
Now define
\begin{align*}
    m(u)
    := \frac{\eta}{e\lambda}e^{\eta(\frac{1}{\lambda}-1)} h(u) - f_{\eta,\lambda}(u).
\end{align*}
By \eqref{eq:appendix_second_derivative_domination}, we have $m''(u) \geq 0$ for all $u>0$, so $m$ is convex.
Also,
\begin{align*}
    h(1) = 0,
    \qquad h'(1) = 0,
    \qquad f_{\eta,\lambda}(1) = 0,
    \qquad f_{\eta,\lambda}'(1) = 0.
\end{align*}
Therefore,
\begin{align*}
    m(1) = 0,
    \qquad m'(1) = 0.
\end{align*}
Since a convex function with derivative $0$ at a point attains its global minimum there, we obtain
\begin{align*}
    m(u) \geq 0
    \qquad \text{for all } u>0.
\end{align*}
This is precisely the pointwise bound \eqref{eq:appendix_pointwise_f_vs_kl}. Integrating
that inequality with respect to $Q$ proves \eqref{eq:appendix_f_vs_kl_divergence}.
\end{proof}

We now evaluate the KL divergence appearing on the right-hand side.
This is exactly the ordinary mutual information between $T_0$ and $X$.

\begin{proposition}[Exact evaluation of the KL mutual information]
\label{prop:appendix_exact_kl_mi}
Assume the relevant differential entropies are finite.
Then
\begin{align}
    \mathrm{KL}(P_{X,T_0}\,\|\,P_XP_{T_0})
    = I(T_0;X)
    = d\log\left(\frac{1}{1-\lambda}\right).
    \label{eq:appendix_exact_kl_mi}
\end{align}
\end{proposition}

\begin{proof}
The identity
\begin{align*}
    \mathrm{KL}(P_{X,T_0}\,\|\,P_XP_{T_0}) = I(T_0;X)
\end{align*}
is the standard definition of mutual information.
So it remains only to compute $I(T_0;X)$.
Using the decomposition
\begin{align*}
    X = \lambda T_0 + (1-\lambda)X',
\end{align*}
with $X'$ independent of $T_0$ and $X' \stackrel{d}{=} X$, we have
\begin{align}
    I(T_0;X)
     & = h(X) - h(X\mid T_0). 
    \label{eq:appendix_mi_entropy_start}
\end{align}
Now condition on $T_0=t$. Then
\begin{align*}
    X\mid(T_0=t) = \lambda t + (1-\lambda)X'.
\end{align*}
The term $\lambda t$ is a deterministic translation, and differential entropy is invariant under translation, hence
\begin{align*}
    h(X\mid T_0=t) = h\left((1-\lambda)X'\right).
\end{align*}
Next, differential entropy scales according to
\begin{align*}
    h(aY) = h(Y) + d\log|a|,
\end{align*}
for every nonzero scalar $a$.
Applying this with $a=1-\lambda$ gives
\begin{align*}
    h\left((1-\lambda)X'\right)
    = h(X') + d\log(1-\lambda).
\end{align*}
Since $X' \stackrel{d}{=} X$, we have $h(X') = h(X)$. Therefore,
\begin{align*}
    h(X\mid T_0=t)
    = h(X) + d\log(1-\lambda),
\end{align*}
and this does not depend on $t$. Averaging over $T_0$ yields
\begin{align*}
    h(X\mid T_0)
    = h(X) + d\log(1-\lambda).
\end{align*}
Substituting this into \eqref{eq:appendix_mi_entropy_start},
\begin{align*}
    I(T_0;X)
     & = h(X) - \left(h(X) + d\log(1-\lambda)\right) \\
     & = d\log\left(\frac{1}{1-\lambda}\right),
\end{align*}
which proves \eqref{eq:appendix_exact_kl_mi}.
\end{proof}

Combining the previous results gives the main estimate.

\begin{corollary}[Explicit KL-based upper bound]
\label{cor:appendix_main_lambda_var_bound}
Under the assumptions above,
\begin{align}
    \Lambda_{\mathrm{var}, q}(E_\lambda, p_X; f)
    \leq
    e^{-\eta} + \eta - 1
    + \frac{\eta}{e\lambda}e^{\eta(\frac{1}{\lambda}-1)}
    d\log\left(\frac{1}{1-\lambda}\right).
    \label{eq:appendix_main_lambda_var_bound}
\end{align}
\end{corollary}

\begin{proof}
Combine \eqref{eq:appendix_f_divergence_bound}, \eqref{eq:appendix_f_vs_kl_divergence}, and \eqref{eq:appendix_exact_kl_mi}.
\end{proof}

The corollary above gives a direct KL-type control of the optimistic-kernel variance term.
It is worth recording separately what happens after the outer factor $1/\eta^2$ is inserted in the final regret bound.

\begin{remark}[After inserting the factor $1/\eta^2$]
\label{rem:appendix_eta_squared}
If the final bound contains $\eta^{-2}\Lambda_{\mathrm{var}, q}(E_\lambda, p_X; f)$, then
\begin{align}
    \frac{1}{\eta^2}\Lambda_{\mathrm{var}, q}(E_\lambda, p_X; f)
    \leq
    \frac{e^{-\eta} + \eta - 1}{\eta^2}
    + \frac{1}{e\lambda\eta}e^{\eta(\frac{1}{\lambda}-1)}
    d\log\left(\frac{1}{1-\lambda}\right).
    \label{eq:appendix_eta_squared_bound}
\end{align}
The first term is harmless for small $\eta$, since
\begin{align*}
    e^{-\eta} + \eta - 1 = \frac{\eta^2}{2} + O(\eta^3),
\end{align*}
so
\begin{align*}
    \frac{e^{-\eta} + \eta - 1}{\eta^2} = \frac{1}{2} + O(\eta).
\end{align*}
Thus the main tuning issue is the second term.
If later one wants to plug in $\eta \asymp 1/\sqrt{T}$, then the present derivation leaves the dependence on
$\lambda$ completely explicit, and one can optimize that tradeoff separately.
I have not tried to optimize that final $\eta$--$\lambda$ tradeoff here; the point of the present subsection is only that the CEXP quantity can be reduced to an ordinary KL mutual information rather than to a $\chi^2$ mutual information.
\end{remark}

Finally, even though the previous corollary already gives a useful upper bound, it is instructive to record that the CEXP quantity stays finite even in the Laplace example where the $\chi^2$ mutual information blows up.

\begin{remark}[The optimistic-kernel CEXP quantity is finite for Laplace]
\label{rem:appendix_laplace_finite}
The quantity in \eqref{eq:appendix_lambda_var_reduce_to_one} is finite for every density $q$ for which the optimistic kernel is defined, and in particular for the Laplace distribution.
Indeed, from \eqref{eq:appendix_phi_after_mean_one} we have
\begin{align*}
    \E_Q\left[\phi_\eta\left(1 - \frac{1}{\lambda} + \frac{1}{\lambda}r\right)\right]
    = \eta - 1 + e^{\eta(\frac{1}{\lambda}-1)}\E_Q\left[e^{-(\eta/\lambda)r}\right].
\end{align*}
Since $r \geq 0$, we have $e^{-(\eta/\lambda)r} \leq 1$, and therefore
\begin{align}
    \Lambda_{\mathrm{var}, q}(E_\lambda, p_X; f)
    \leq
    \eta - 1 + e^{\eta(\frac{1}{\lambda}-1)} < \infty.
    \label{eq:appendix_uniform_finiteness}
\end{align}
So the Laplace pathology that appears for the $\chi^2$ mutual information does \emph{not} appear here.
The structural reason is that the $\chi^2$ generator grows quadratically, whereas the generator
$f_{\eta,\lambda}$ in \eqref{eq:appendix_f_generator_def} has only linear growth as $u \to \infty$.
\end{remark}
